\documentclass[10pt]{article}
\usepackage[portuguese]{babel}
\usepackage[a4paper, top = 0.8in, bottom = 0.9in, left =0.6in, right = 0.6in]{geometry}
\usepackage{inputenc}
\usepackage{fancyhdr}
\usepackage{multirow}
\usepackage{graphicx}
\usepackage{times}
\usepackage{subcaption}
\usepackage{booktabs}
\usepackage{microtype}
\usepackage{wrapfig}
\usepackage{circuitikz}
\usepackage{titlesec}
\usepackage{wrapfig}
\usepackage{listings}
\usepackage{tikz}

\usepackage{caption}
\usepackage{ragged2e}
\usepackage{adjustbox}
\usepackage{float}
\usepackage{amsmath, mathtools, amssymb}
\usepackage[makeroom]{cancel}
\usepackage{parskip}
\usepackage{array, makecell}
\usepackage{booktabs}
\titleformat{\section}
{\normalfont\large\bfseries}{\thesection)}{1em}{}
\titleformat{\subsection}
{\normalfont\normalsize\bfseries}{\thesubsection)}{1em}{}
\titleformat{\subsubsection}
{\normalfont\normalsize\bfseries}{\thesubsubsection)}{1em}{}


%\pretitle{\begin{center}\Huge\bfseries}

\pagestyle{fancy}
\fancyhead{}
\centering\chead{\includegraphics[width=15cm]{imagens/unb_fga_logo_extenso.jpg}}
\lfoot{Prática de Circuitos 1 - 2020/1}
\rfoot{\textbf{Prof. Dr. Marcus V. Chaffim Costa}} %nome do rodape
\renewcommand{\footrulewidth}{1pt}

\setlength{\parskip}{0.1cm}
\author{Monitoria - 2019/2}
\title{Tutorial 0}


%arrumação das tabelas

\setlength{\tabcolsep}{0.5em} % for the horizontal padding
{\renewcommand{\arraystretch}{1.2}}% for the vertical padding
\captionsetup[table]{name= \bfseries Tabela}
\captionsetup[figure]{name= \bfseries Figura}
\numberwithin{table}{section}
\numberwithin{figure}{section}
\usepackage{array}
\newcolumntype{C}[1]{>{\centering\arraybackslash}m{#1}}

\aboverulesep=0cm
\belowrulesep=0cm

%configuracoes dos codigos em matlab

\usepackage{color} %red, green, blue, yellow, cyan, magenta, black, white
\definecolor{mygreen}{RGB}{28,172,0} % color values Red, Green, Blue
\definecolor{mylilas}{RGB}{170,55,241}

\lstset{language=Matlab,%
    %basicstyle=\color{red},
    breaklines=true,%
    morekeywords={matlab2tikz},
    keywordstyle=\color{blue},%
    morekeywords=[2]{1}, keywordstyle=[2]{\color{black}},
    identifierstyle=\color{black},%
    stringstyle=\color{mylilas},
    commentstyle=\color{mygreen},%
    showstringspaces=false,%without this there will be a symbol in the places where there is a space
    numbers=left,%
    numberstyle={\tiny \color{black}},% size of the numbers
    numbersep=9pt, % this defines how far the numbers are from the text
    emph=[1]{for,end,break},emphstyle=[1]\color{red}, %some words to emphasise
    %emph=[2]{word1,word2}, emphstyle=[2]{style},    
}

\allowdisplaybreaks


%==============macros para uso nos documentos===================%

\newcommand{\ohm}{\Omega}


\newcommand{\figuras}[4]{
    \begin{figure}[H]
        \centering
        \includegraphics[width={#1}\textwidth]{{#2}}
        \caption{{#3}}
        \label{fig:{#4}}
    \end{figure}
}

\newcommand{\figura}[2]{
    \begin{figure}[H]
        \centering
        \includegraphics[width={#1}\textwidth]{{#2}}
    \end{figure}
}

\newcommand{\subfiguras}[5]{
    \centering
    \begin{subfigure}[H]{#1\textwidth}
        \centering
        \includegraphics[width={#2}\textwidth]{{#3}}
        \caption{{#4}}
        \label{fig:{#5}}
    \end{subfigure}
}

\newcommand{\circuito}[3]{
    \begin{figure}[H]
        \centering
        \begin{circuitikz}[line width=.5pt]
            #1
        \end{circuitikz}
        \caption{{#2}}
        \label{circ:{#3}}    
    \end{figure}
}

\newcommand{\calc}[1]{
    \begin{gather*}
        #1
    \end{gather*}
}





%=====================================================%

\usetikzlibrary{automata}

\begin{document}

\begin{center}
\vspace*{.03cm}
\Large\textbf{Experimento 02:}\\ %titulo
\Large{Leis de Kirchhoff}
\end{center}
\justify

\section{Objetivos}

O objetivo deste experimento é a prática do método dos nós e das malhas para análise de circuitos eletrônicos. São abordadas
as resoluções de circuitos utilizando estes métodos, bem como a montagem dos circuitos correspondentes e sua verificação
experimental através da medição das grandezas elétricas consideradas nas resoluções.

\section{Objetivos}
Nesta experiência investigaremos alguns dos instrumentos de bancada do laboratório a fim de compreender como
funcionam. Em seguida produziremos sinais DC e realizaremos medições de tensão,corrente e resistência. O intuito é desenvolver habilidades na manipulação dos equipamentos.

\section{Estudo pré-laboratorial}

\subsection{Instrumentos de Bancada}

\subsubsection{Fonte de alimentação}

\begin{itemize}
    \item \textbf{O que é faixa de tensão de saída do equipamento?}
        
        É a faixa de tensões em que o aparelho fornece em sua saída. Ela pode ser limitada nominalmente pelo usuário através de um potenciômetro (em algumas literaturas, \textit{knob}). A faixa de tensões vária em função do modo de operação escolhido entre os canais de saída da fonte.

    \item \textbf{O que é faixa de corrente de saída do equipamento?}

        É a faixa de corrente que o aparelho fornece em sua saída. É análogo à faixa de tensões, mas conta, vale ressaltar, com o ajuste de corrente limite. A faixa de corrente também vária em função do modo de operação configurada na fonte.
    \item \textbf{O que é corrente limite?}

        É a corrente máxima configurada pelo usuário a partir da qual a proteção contra sobrecorrente atua. A quantidade de corrente opera a partir da resistência oferecida pelo circuito, ou seja, se há algum componente que esteja danificado, há probabilidades que haja o sobreaquecimento e, com isso, até a combustão do componente.
\end{itemize}

\subsubsection{Modos de Operação de Fontes de Alimentação}

Atualmente há, majoritariamente, dois tipos de fonte no laboratório: MPL-1303 e MPL-3305, ambas da Minipa e, portanto, compartilham do mesmo manual. A principal diferença reside entre a quantidade de canais de saída existentes: a MPL-1303 possui apenas um canal de saída, enquanto a MPL-3305 possui dois canais de saída. Implica-se, então, que a MPL-1303 possui apenas o modo de operação simples.

Para todas os modos de operação, devemos configurar as teclas de Seleção do Modo de Conexão (também conhecidas como Teclas de Tracking). Elas mudam a conexão entre as fontes de acordo com as associações desejadas. As configurações de Tracking dos modos Independente, Série ou paralelo para a fonte MPL-3305 estão elencados abaixo:

\figuras{.5}{imagens/instr_banc/tracking}{Configurações de Conexão}{tracking}

\begin{itemize}
     
    \item \textbf{Fixa}

        A fonte MPL-3305 conta com três saídas, sendo duas variáveis e uma fixa de 5V.
        
        \newpage

    \item \textbf{Simples}

        É o modo onde os canais de saída operam independentemente. A corrente máxima é de 5A.

        \figuras{.5}{imagens/instr_banc/op_simples}{Modo de Operação Simples.}{op_simples}
    \item \textbf{Paralelo}

        Nessa condição, os canais operam paralelamente e, consequentemente, as corrente se somam. A corrente máxima é, portanto, 10A.

        \figuras{.4}{imagens/instr_banc/op_paralelo}{Modo de Operação Paralela.}{op_paralelo}

    \item \textbf{Série}

        Nessa condição, os canais operam em serialmente e, consequentemente, as tensões se somam. A tensão máxima é 64V.

        \figuras{.4}{imagens/instr_banc/op_serie}{Modo de Operação Série.}{op_serie}
    \item \textbf{Simétrica}
    
        Nesta condição, pode-se conseguir um terra comum para ambos os canais variáveis, com saída positiva e negativa de, no máximo, +32V e -32V, respectivamente. 
    
        \figuras{.4}{imagens/instr_banc/op_simetric}{Modo de Operação Simétrica.}{op_simetrica}
\end{itemize}

O ajuste de corrente limite é ilustrado na figura abaixo:

\figuras{.4}{imagens/instr_banc/corr_lim}{Configuração da Corrente Limite na fonte MPL-3305}{}

Observe que, propositalmente, não houve um detalhamento das peculiaridades de cada um dos modos de operação. Através dos manuais dos equipamentos, procure compreender o modo de funcionamento das fontes de alimentação. Como guia, responda as seguintes perguntas:

a) Fonte de alimentação modelo MPL-1303 Minipa:

\begin{itemize}
    \item Qual a faixa de tensão de saída desse equipamento?
    \item Qual a faixa de corrente de saída desse equipamento?
    \item Como definir um valor máximo de corrente de saída para a fonte (corrente limite)?
\end{itemize}

\vspace{.5em}

b) Fonte de alimentação modelo MPL-3305 Minipa:

\begin{itemize}
    \item Qual a faixa de tensão de saída para esse equipamento?
    \item Qual a faixa de corrente de saída para esse equipamento?
    \item Quais os limites de tensão e corrente obtidos em cada um dos modos de operação?
\end{itemize}

\subsubsection{Multímetro}

\begin{itemize}
    \item \textbf{Medição de Tensão com Multímetro:}Posicione a chave rotativa para medição de voltagem contínua e conecte as pontas de prova em paralelo no circuito em teste.
    
    \figuras{.2}{imagens/instr_banc/mult_tens}{Medição de tensão com o multímetro.}{tens_mult}

    \item \textbf{Medição de corrente com o Multímetro:} Posicione a chave rotativa na maior escala possível de corrente e ajuste para uma visualização adequada. Conecte as pontas de prova em série no local a ser medido.
    
    \figuras{.35}{imagens/instr_banc/corr_mult}{Medição de corrente com o multímetro.}{corr_mult}
    
    \item \textbf{Medição de Resistência com o multímetro:} Posicione a chave rotativa em "$\ohm$". Conecte as pontas de prova sobre o objeto a ser medido.
    
    \figuras{.2}{imagens/instr_banc/res_mult}{Medição de Resistência com o multímetro.}{res_mult}

    \item \textbf{Integridade de Trilhas:} Posicione a chave rotativa para medição de continuidade. O aparelho sonoriza um \textit{beep} quando a resistência de um circuito for menor que $10\ohm$.
    
    \figuras{.2}{imagens/instr_banc/int_trilh}{Teste de integridade de trilhas com o multímetro.}{int_trilh}
\end{itemize}

Assim como as fontes de tensão, há dois modelos de multímetro disponíveis no laboratório atualmente: ET-1110 da Minipa e o TOOZ DT830 Series. Estude a manipulação e as características nos manuais do aparelho.





\end{document}